% Options for packages loaded elsewhere
% Options for packages loaded elsewhere
\PassOptionsToPackage{unicode}{hyperref}
\PassOptionsToPackage{hyphens}{url}
%
\documentclass[
  ignorenonframetext,
  aspectratio=169,
]{beamer}
\newif\ifbibliography
\usepackage{pgfpages}
\setbeamertemplate{caption}[numbered]
\setbeamertemplate{caption label separator}{: }
\setbeamercolor{caption name}{fg=normal text.fg}
\beamertemplatenavigationsymbolsempty
% remove section numbering
\setbeamertemplate{part page}{
  \centering
  \begin{beamercolorbox}[sep=16pt,center]{part title}
    \usebeamerfont{part title}\insertpart\par
  \end{beamercolorbox}
}
\setbeamertemplate{section page}{
  \centering
  \begin{beamercolorbox}[sep=12pt,center]{section title}
    \usebeamerfont{section title}\insertsection\par
  \end{beamercolorbox}
}
\setbeamertemplate{subsection page}{
  \centering
  \begin{beamercolorbox}[sep=8pt,center]{subsection title}
    \usebeamerfont{subsection title}\insertsubsection\par
  \end{beamercolorbox}
}
% Prevent slide breaks in the middle of a paragraph
\widowpenalties 1 10000
\raggedbottom
\AtBeginPart{
  \frame{\partpage}
}
\AtBeginSection{
  \ifbibliography
  \else
    \frame{\sectionpage}
  \fi
}
\AtBeginSubsection{
  \frame{\subsectionpage}
}
\usepackage{iftex}
\ifPDFTeX
  \usepackage[T1]{fontenc}
  \usepackage[utf8]{inputenc}
  \usepackage{textcomp} % provide euro and other symbols
\else % if luatex or xetex
  \usepackage{unicode-math} % this also loads fontspec
  \defaultfontfeatures{Scale=MatchLowercase}
  \defaultfontfeatures[\rmfamily]{Ligatures=TeX,Scale=1}
\fi
\usepackage{lmodern}

\ifPDFTeX\else
  % xetex/luatex font selection
\fi
% Use upquote if available, for straight quotes in verbatim environments
\IfFileExists{upquote.sty}{\usepackage{upquote}}{}
\IfFileExists{microtype.sty}{% use microtype if available
  \usepackage[]{microtype}
  \UseMicrotypeSet[protrusion]{basicmath} % disable protrusion for tt fonts
}{}
\makeatletter
\@ifundefined{KOMAClassName}{% if non-KOMA class
  \IfFileExists{parskip.sty}{%
    \usepackage{parskip}
  }{% else
    \setlength{\parindent}{0pt}
    \setlength{\parskip}{6pt plus 2pt minus 1pt}}
}{% if KOMA class
  \KOMAoptions{parskip=half}}
\makeatother


\usepackage{longtable,booktabs,array}
\usepackage{calc} % for calculating minipage widths
\usepackage{caption}
% Make caption package work with longtable
\makeatletter
\def\fnum@table{\tablename~\thetable}
\makeatother
\usepackage{graphicx}
\makeatletter
\newsavebox\pandoc@box
\newcommand*\pandocbounded[1]{% scales image to fit in text height/width
  \sbox\pandoc@box{#1}%
  \Gscale@div\@tempa{\textheight}{\dimexpr\ht\pandoc@box+\dp\pandoc@box\relax}%
  \Gscale@div\@tempb{\linewidth}{\wd\pandoc@box}%
  \ifdim\@tempb\p@<\@tempa\p@\let\@tempa\@tempb\fi% select the smaller of both
  \ifdim\@tempa\p@<\p@\scalebox{\@tempa}{\usebox\pandoc@box}%
  \else\usebox{\pandoc@box}%
  \fi%
}
% Set default figure placement to htbp
\def\fps@figure{htbp}
\makeatother


% definitions for citeproc citations
\NewDocumentCommand\citeproctext{}{}
\NewDocumentCommand\citeproc{mm}{%
  \begingroup\def\citeproctext{#2}\cite{#1}\endgroup}
\makeatletter
 % allow citations to break across lines
 \let\@cite@ofmt\@firstofone
 % avoid brackets around text for \cite:
 \def\@biblabel#1{}
 \def\@cite#1#2{{#1\if@tempswa , #2\fi}}
\makeatother
\newlength{\cslhangindent}
\setlength{\cslhangindent}{1.5em}
\newlength{\csllabelwidth}
\setlength{\csllabelwidth}{3em}
\newenvironment{CSLReferences}[2] % #1 hanging-indent, #2 entry-spacing
 {\begin{list}{}{%
  \setlength{\itemindent}{0pt}
  \setlength{\leftmargin}{0pt}
  \setlength{\parsep}{0pt}
  % turn on hanging indent if param 1 is 1
  \ifodd #1
   \setlength{\leftmargin}{\cslhangindent}
   \setlength{\itemindent}{-1\cslhangindent}
  \fi
  % set entry spacing
  \setlength{\itemsep}{#2\baselineskip}}}
 {\end{list}}
\usepackage{calc}
\newcommand{\CSLBlock}[1]{\hfill\break\parbox[t]{\linewidth}{\strut\ignorespaces#1\strut}}
\newcommand{\CSLLeftMargin}[1]{\parbox[t]{\csllabelwidth}{\strut#1\strut}}
\newcommand{\CSLRightInline}[1]{\parbox[t]{\linewidth - \csllabelwidth}{\strut#1\strut}}
\newcommand{\CSLIndent}[1]{\hspace{\cslhangindent}#1}



\setlength{\emergencystretch}{3em} % prevent overfull lines

\providecommand{\tightlist}{%
  \setlength{\itemsep}{0pt}\setlength{\parskip}{0pt}}



 


\makeatletter
\@ifpackageloaded{caption}{}{\usepackage{caption}}
\AtBeginDocument{%
\ifdefined\contentsname
  \renewcommand*\contentsname{Table of contents}
\else
  \newcommand\contentsname{Table of contents}
\fi
\ifdefined\listfigurename
  \renewcommand*\listfigurename{List of Figures}
\else
  \newcommand\listfigurename{List of Figures}
\fi
\ifdefined\listtablename
  \renewcommand*\listtablename{List of Tables}
\else
  \newcommand\listtablename{List of Tables}
\fi
\ifdefined\figurename
  \renewcommand*\figurename{Figure}
\else
  \newcommand\figurename{Figure}
\fi
\ifdefined\tablename
  \renewcommand*\tablename{Table}
\else
  \newcommand\tablename{Table}
\fi
}
\@ifpackageloaded{float}{}{\usepackage{float}}
\floatstyle{ruled}
\@ifundefined{c@chapter}{\newfloat{codelisting}{h}{lop}}{\newfloat{codelisting}{h}{lop}[chapter]}
\floatname{codelisting}{Listing}
\newcommand*\listoflistings{\listof{codelisting}{List of Listings}}
\makeatother
\makeatletter
\makeatother
\makeatletter
\@ifpackageloaded{caption}{}{\usepackage{caption}}
\@ifpackageloaded{subcaption}{}{\usepackage{subcaption}}
\makeatother

\usepackage{bookmark}
\IfFileExists{xurl.sty}{\usepackage{xurl}}{} % add URL line breaks if available
\urlstyle{same}
\hypersetup{
  pdftitle={Parte 1 : Regressão},
  hidelinks,
  pdfcreator={LaTeX via pandoc}}



\title{Parte 1 : Regressão}
\subtitle{{Métodos Paramétricos}}
\author{}
\date{}

\begin{document}
\frame{\titlepage}


\begin{frame}{Métodos paramétricos}
\phantomsection\label{muxe9todos-paramuxe9tricos}
Métodos \textbf{paramétricos} assumem que a função de regressão pode ser
descrita por um \textbf{número finito de parâmetros}.

Aqui, iremos nos concentrar em \textbf{modelos de regressão linear}.

\begin{itemize}
\tightlist
\item
  Grande parte do conteúdo pode parecer familiar;
\item
  porém, será apresentado sob a perspectiva de \textbf{modelagem
  preditiva e seleção de modelos}.
\end{itemize}
\end{frame}

\begin{frame}{Regressão linear}
\phantomsection\label{regressuxe3o-linear}
Na regressão linear, a função de predição possui \textbf{forma linear
nos parâmetros}.

A função de predição pode ser escrita como

\[
g(x) = \beta^t x
\]

ou explicitamente

\[
g(x) = \beta_0 x_0 + \beta_1 x_1 + \cdots + \beta_p x_p
\] Adotamos a convenção \(x_0 = 1\). Assim, \(\beta_0\) corresponde ao
\textbf{intercepto do modelo}.

O vetor de parâmetros é

\[
\beta = (\beta_0, \beta_1, \ldots, \beta_p)
\]
\end{frame}

\begin{frame}{Transformações de variáveis}
\phantomsection\label{transformauxe7uxf5es-de-variuxe1veis}
Os termos \(x_i\) \textbf{não precisam ser apenas as variáveis
originais}. Podemos construir \textbf{novas covariáveis} a partir delas,
por exemplo:

\begin{itemize}
\tightlist
\item
  termos quadráticos: \(x_i^2\)
\item
  interações: \(x_i x_j\)
\item
  transformações não lineares
\end{itemize}

Isso permite que \textbf{modelos lineares representem relações não
lineares nos dados}.
\end{frame}

\begin{frame}{O Método dos Mínimos Quadrados}
\phantomsection\label{o-muxe9todo-dos-muxednimos-quadrados}
Uma forma de estimar os coeficientes \(\beta\) da regressão linear é
utilizar o \textbf{método dos mínimos quadrados}.

A ideia é escolher os coeficientes que \textbf{minimizam a soma dos
erros quadráticos}:

\[
\hat{\beta} =
\arg\min_{\beta}
\sum_{i=1}^{n}
\left(
Y_i - \beta_0 - \beta_1 x_{i1} - \cdots - \beta_d x_{ip}
\right)^2
\]
\end{frame}

\begin{frame}{Interpretação}
\phantomsection\label{interpretauxe7uxe3o}
O método escolhe os parâmetros que \textbf{minimizam a soma dos
quadrados dos resíduos}:

\[
\text{resíduo}_i = Y_i - g(X_i)
\]

Logo, o objetivo é minimizar

\[
\sum_{i=1}^{n} (Y_i - g(X_i))^2
\]

Geometricamente, isso corresponde a \textbf{ajustar o hiperplano que
fica mais próximo possível dos dados}.
\end{frame}

\begin{frame}{Forma matricial do modelo}
\phantomsection\label{forma-matricial-do-modelo}
O modelo de regressão pode ser escrito em forma matricial como

\[
\boldsymbol{y} = \boldsymbol{X} \boldsymbol{\beta} + \boldsymbol{\varepsilon}
\]

onde

\begin{itemize}
\tightlist
\item
  \(\boldsymbol{y}\) é o vetor de respostas
\item
  \(\boldsymbol{X}\) é a \textbf{matriz de covariáveis}
\item
  \(\boldsymbol{\beta}\) é o vetor de parâmetros
\end{itemize}
\end{frame}

\begin{frame}{Solução dos mínimos quadrados}
\phantomsection\label{soluuxe7uxe3o-dos-muxednimos-quadrados}
A solução analítica do problema é dada por

\[
\hat{\boldsymbol{\beta}} = (\boldsymbol{X}^t \boldsymbol{X})^{-1} \boldsymbol{X}^t \boldsymbol{y}
\]

Essa expressão é conhecida como a \textbf{equação normal} da regressão
linear.
\end{frame}

\begin{frame}{Estrutura da matriz de dados}
\phantomsection\label{estrutura-da-matriz-de-dados}
A matriz de covariáveis possui a forma

\[
\boldsymbol{X} =
\begin{pmatrix}
x_{1,0} & x_{1,1} & \cdots & x_{1,d} \\
x_{2,0} & x_{2,1} & \cdots & x_{2,d} \\
\vdots  & \vdots  & \ddots & \vdots \\
x_{n,0} & x_{n,1} & \cdots & x_{n,d}
\end{pmatrix}
\]
\end{frame}

\begin{frame}
e o vetor resposta é

\[
\boldsymbol{y} =
\begin{pmatrix}
Y_1 \\
Y_2 \\
\vdots \\
Y_n
\end{pmatrix}
\]

Adotamos a convenção \(x_{i,0} = 1\) para representar o
\textbf{intercepto} do modelo.
\end{frame}

\begin{frame}{Função de regressão estimada}
\phantomsection\label{funuxe7uxe3o-de-regressuxe3o-estimada}
Uma vez estimados os coeficientes pelo método dos mínimos quadrados, a
\textbf{função de predição} é dada por

\[
g(\mathbf{x}) = \hat{\boldsymbol{\beta}}^t \mathbf{x}
\]

Assim, para um novo vetor de covariáveis \(\mathbf{x}\), a previsão é

\[
\hat{\boldsymbol{y}} = g(\mathbf{x})
\]
\end{frame}

\begin{frame}{Justificativa estatística do método}
\phantomsection\label{justificativa-estatuxedstica-do-muxe9todo}
Grande parte da literatura estatística estuda o método dos mínimos
quadrados sob diferentes perspectivas, como:

\begin{itemize}
\tightlist
\item
  \textbf{máxima verossimilhança};
\item
  \textbf{testes de hipóteses} para os coeficientes \(\beta_i\);
\item
  \textbf{intervalos de confiança} para os parâmetros.
\end{itemize}

Esses tópicos são tratados em profundidade em livros clássicos como:

\begin{itemize}
\tightlist
\item
  Morettin and Singer (2019)
\item
  Neter et al. (1996)
\end{itemize}
\end{frame}

\begin{frame}{Uma observação importante}
\phantomsection\label{uma-observauxe7uxe3o-importante}
Assumir que a verdadeira função de regressão \(r(x)\) é exatamente
linear é uma \textbf{suposição bastante forte}.

Na prática, essa hipótese raramente é verdadeira.

Mesmo assim, existe uma literatura extensa mostrando que o método dos
mínimos quadrados continua sendo útil para estimar relações entre
\(\boldsymbol{X}\) e \(\boldsymbol{y}\), mesmo quando a verdadeira
função de regressão \textbf{não é linear}.
\end{frame}

\begin{frame}{Quando a suposição de linearidade falha}
\phantomsection\label{quando-a-suposiuxe7uxe3o-de-linearidade-falha}
A suposição de que a verdadeira função de regressão

\[
r(\mathbf{x}) = E[\boldsymbol{y} \mid \boldsymbol{X}=\mathbf{x}]
\]

é \textbf{linear} muitas vezes \textbf{não é válida na prática}.

Mesmo assim, modelos lineares frequentemente apresentam \textbf{bom
desempenho preditivo}.
\end{frame}

\begin{frame}
\begin{block}{Ideia central}
\phantomsection\label{ideia-central}
Mesmo quando \(r(\mathbf{x})\) não é linear, frequentemente existe um
vetor \(\boldsymbol{\beta}_*\) tal que o preditor linear

\[
g_{\boldsymbol{\beta}_*}(\mathbf{x}) = \boldsymbol{\beta}_*^t \mathbf{x}
\]

possui \textbf{bom poder preditivo}.
\end{block}
\end{frame}

\begin{frame}
\begin{block}{O melhor preditor linear}
\phantomsection\label{o-melhor-preditor-linear}
O vetor \(\boldsymbol{\beta}_*\) é definido como aquele que
\textbf{minimiza o risco quadrático}:

\[
\boldsymbol{\beta}_* =
\arg\min_{\boldsymbol{\beta}}
R(g_{\boldsymbol{\beta}})
=
\arg\min_{\boldsymbol{\beta}}
E[(\boldsymbol{y} - \boldsymbol{\beta}^t \boldsymbol{X})^2]
\]

Esse vetor define a \textbf{melhor aproximação linear} da relação entre
\(\boldsymbol{X}\) e \(\boldsymbol{y}\).

\(\boldsymbol{\beta}_*\) é chamado de \textbf{oráculo}.
\end{block}
\end{frame}

\begin{frame}
\begin{block}{Interpretação}
\phantomsection\label{interpretauxe7uxe3o-1}
Mesmo quando o verdadeiro modelo \(r(x)\) não é linear, o método dos
mínimos quadrados estima o \textbf{melhor preditor linear possível}.

Ou seja, ele encontra a função linear que \textbf{minimiza o erro
quadrático médio esperado}.
\end{block}
\end{frame}

\begin{frame}
\begin{block}{Consequência importante}
\phantomsection\label{consequuxeancia-importante}
Quando o tamanho da amostra cresce,

\[
\hat{\boldsymbol{\beta}} \rightarrow \boldsymbol{\beta}_* 
\]

Ou seja, o estimador obtido pelos \textbf{mínimos quadrados converge
para o melhor preditor linear}, mesmo que o modelo esteja \textbf{mal
especificado}.

\begin{quote}
\textbf{Mínimos quadrados não precisa que o modelo seja verdadeiro.}\\
Ele encontra a \textbf{melhor aproximação linear da realidade}.
\end{quote}
\end{block}
\end{frame}

\begin{frame}{Limitações do método dos mínimos quadrados}
\phantomsection\label{limitauxe7uxf5es-do-muxe9todo-dos-muxednimos-quadrados}
Quando o número de covariáveis é \textbf{muito grande}, o método dos
mínimos quadrados pode apresentar \textbf{desempenho ruim}.

Isso ocorre principalmente por dois motivos:

\begin{itemize}
\tightlist
\item
  \textbf{super-ajuste (overfitting)}\\
\item
  \textbf{alta variância do estimador}
\end{itemize}
\end{frame}

\begin{frame}
\begin{block}{Problema em alta dimensão}
\phantomsection\label{problema-em-alta-dimensuxe3o}
Considere:

\begin{itemize}
\tightlist
\item
  \(p\) = número de covariáveis\\
\item
  \(n\) = tamanho da amostra
\end{itemize}

Quando \(p > n\) o estimador de mínimos quadrados \textbf{não está bem
definido}.
\end{block}
\end{frame}

\begin{frame}
\begin{block}{Por quê?}
\phantomsection\label{por-quuxea}
A solução dos mínimos quadrados é

\[
\hat{\boldsymbol{\beta}} = (\boldsymbol{X}^t \boldsymbol{X})^{-1} \boldsymbol{X}^t \boldsymbol{y}
\]

Porém, quando \(p > n\), \(\boldsymbol{X}^t \boldsymbol{X}\) \textbf{não
é invertível}.

Logo, o estimador não pode ser calculado diretamente.
\end{block}
\end{frame}

\begin{frame}
\begin{block}{Consequência}
\phantomsection\label{consequuxeancia}
Em problemas com \textbf{muitas variáveis} (alta dimensão):

\begin{itemize}
\tightlist
\item
  o método clássico de mínimos quadrados \textbf{não funciona bem}
\item
  precisamos utilizar \textbf{métodos alternativos}
\end{itemize}

Por exemplo:

\begin{itemize}
\tightlist
\item
  \textbf{regularização (Ridge, Lasso)}
\item
  \textbf{redução de dimensionalidade}
\item
  \textbf{seleção de variáveis}
\end{itemize}
\end{block}
\end{frame}

\begin{frame}
\pandocbounded{\includegraphics[keepaspectratio]{met_param_files/figure-beamer/unnamed-chunk-1-1.png}}
\end{frame}

\begin{frame}
\begin{block}{Interpretação}
\phantomsection\label{interpretauxe7uxe3o-2}
\begin{quote}
Quando temos mais parâmetros do que observações, o modelo tem liberdade
suficiente para passar exatamente por todos os pontos.
\end{quote}

\begin{quote}
O problema é que existem infinitas funções que fazem isso, e elas podem
ser completamente diferentes fora da amostra.
\end{quote}
\end{block}
\end{frame}

\begin{frame}{References}
\phantomsection\label{references}
\phantomsection\label{refs}
\begin{CSLReferences}{1}{0}
\bibitem[\citeproctext]{ref-morettin2019estatistica}
Morettin, Pedro A., and Julio M. Singer. 2019. \emph{Estatística e
Ciência de Dados}. Rio de Janeiro: LTC.

\bibitem[\citeproctext]{ref-neter1996applied}
Neter, John, Michael H. Kutner, Christopher J. Nachtsheim, and William
Wasserman. 1996. \emph{Applied Linear Statistical Models}. 4th ed.
Chicago: Irwin.

\end{CSLReferences}
\end{frame}




\end{document}
